\chapter{Introduction}
\label{ch:intro}

\section{Background and motivation}
When people choose a restaurant they rarely think in categories. They think in situations and
tastes: ``I am on a budget'', ``a bar with a patio'', or, in Thai, \thai{อยากกินอะไรเปรี้ยวๆ}
(``I want something sour''). Most restaurant platforms answer such needs with filters, star
ratings and popularity. These signals are easy to compute, but they do not say whether a
restaurant fits \emph{this} request. A five-star seafood restaurant is a poor answer to a
student on a budget.

The information needed to answer these requests usually exists already. Diners describe
prices, waiting times, noise, portion sizes and flavours in their reviews. A system that
searches review text directly can match a free-text request to the restaurants whose
reviewers describe what the user wants. This task is called \emph{reviewed-item retrieval}
\citep{rird2023}: items are represented only by their reviews, and the system ranks items for a
natural-language query.

Searching reviews has a second advantage: it produces its own evidence. If a restaurant is
ranked first because of a review sentence, that sentence can be shown to the user as the
reason. The user can then read it and decide whether to trust the recommendation. This
evidence is selected from real reviews, not written by a language model, so it cannot invent
facts about the restaurant. The evidence is useful only if it is correct, and
Figure~\ref{fig:intro-example} shows a case where it is not. For the query ``I am on a
budget'', our proposed model ranks a restaurant that annotators judged relevant first, with a
sentence that clearly supports the query. The second result is a restaurant that all five
annotators judged \emph{not} relevant, and the sentence shown for it, ``Be prepared to spend
\$\$'', says the opposite of what the user asked for. The model matched the topic (price) but
not the meaning (cheap). A ranking with evidence is only trustworthy if we measure whether the
evidence actually supports the ranking, and that is one of the two central questions of this
report.

\begin{figure}[H]
\centering
\begin{tikzpicture}[font=\small,
  box/.style={draw,rounded corners,inner sep=5pt},
  arrow/.style={-{Latex[length=2mm]},thick}]
\node[box,fill=blue!7,text width=3.0cm,align=center] (q) {\textbf{User query}\\``I am on a budget''};
\node[box,fill=gray!10,text width=3.0cm,align=center,below=9mm of q] (s)
  {Dense sentence retrieval,\\top-10 mean};
\node[box,fill=green!8,text width=8.4cm,align=left,anchor=north west] (r1)
  at ([xshift=12mm]q.north east)
  {\textbf{\#1 Manpuku Japanese Eatery} \hfill relevant (5 of 5 annotators)\\[2pt]
   \emph{``Again, not amazing but will do if you are on a budget.''}};
\node[box,fill=red!7,text width=8.4cm,align=left,below=5mm of r1] (r2)
  {\textbf{\#2 Fishman Lobster Clubhouse} \hfill not relevant (0 of 5)\\[2pt]
   \emph{``Be prepared to spend \$\$''}};
\draw[arrow] (q) -- (s);
\draw[arrow] (s.east) -- (r1.west);
\draw[arrow] (s.east) -- (r2.west);
\end{tikzpicture}
\caption{A real output of the proposed model on RIRD. Each restaurant is shown with its
highest-scoring review sentence. The first piece of evidence supports the query; the second
matches its topic but contradicts it. Relevance labels are the RIRD annotator votes
\citep{rird2023}; the output is generated by this project.}
\label{fig:intro-example}
\end{figure}

\section{From the proposal to the final scope}
The project started as a proposal for a Thai food recommender that would match taste
preferences such as sour, spicy or not too sweet. A data audit and the feedback on the
proposal changed this plan. The audit found that no available dataset links Thai reviews to
restaurants and also contains relevance labels for preference queries. The English
Reviewed-Item Retrieval Dataset (RIRD) has restaurant identifiers and human relevance
judgements for 100 queries \citep{rird2023}, but its reviews are about Toronto restaurants and
its queries are general preferences, not Thai tastes. The Thai Wongnai dataset has real Thai
reviews \citep{wongnai2020}, but no restaurant identifiers and no query labels.

The proposal feedback raised four concerns. First, the application claims (Thai taste) did not
match what was evaluated (English restaurant preferences). Second, BM25 and dense retrieval
were compared on different text units, and restaurants with more reviews could be favoured.
Third, the quality of the shown explanations had not been tested. Fourth, more literature was
needed. The feedback ended with a question that we adopted as the guiding question of the final
project:
\begin{quote}
\emph{Do retrieved review sentences genuinely support the system's restaurant rankings, and can
the same approach work on real Thai restaurant reviews with independently annotated taste
preferences?}
\end{quote}

We therefore split the work into two tracks and limited each claim to what its data can
measure (Figure~\ref{fig:intro-scope}). The English track measures the full system: restaurant
ranking and the explanations shown for each restaurant. The Thai track measures only the first
step, finding review sentences that match a taste query, on a new human-labelled Thai test set
that we built for this purpose. Thai restaurant ranking is left as future work because it
cannot be evaluated without restaurant identifiers.

\begin{figure}[H]
\centering
\begin{tikzpicture}[font=\small,
  blk/.style={draw,rounded corners,align=center,minimum height=11mm,inner sep=4pt},
  arrow/.style={-{Latex[length=2mm]},thick}]
\node[blk,fill=orange!12,text width=2.6cm] (p) {\textbf{Proposal}\\Thai food\\recommender};
\node[blk,fill=gray!10,text width=2.6cm,right=8mm of p] (a) {Data audit and\\proposal feedback};
\node[blk,fill=blue!8,text width=5.0cm,right=12mm of a,yshift=16mm] (en)
  {\textbf{English track (RIRD)}\\restaurant ranking and explanations,\\measured with existing labels};
\node[blk,fill=green!8,text width=5.0cm,right=12mm of a] (th)
  {\textbf{Thai track (Wongnai)}\\sentence retrieval for taste queries,\\measured with new human labels};
\node[blk,dashed,fill=gray!5,text width=5.0cm,right=12mm of a,yshift=-16mm] (fw)
  {\textbf{Thai restaurant ranking}\\future work: no restaurant IDs};
\draw[arrow] (p) -- (a);
\draw[arrow] (a.east) -- (en.west);
\draw[arrow] (a.east) -- (th.west);
\draw[arrow,dashed,gray] (a.east) -- (fw.west);
\end{tikzpicture}
\caption{How the data audit and the proposal feedback turned one broad claim into two
measurable tracks.}
\label{fig:intro-scope}
\end{figure}

\section{Business understanding and potential impact}
The intended users are diners who choose a restaurant from a personal need rather than a fixed
category. The business problem is not ``predict the highest-rated restaurant''. A user may
prefer a lower-rated restaurant because it is cheap, quiet, quick or suitable for children. A
single global rating cannot express these trade-offs, while review text often can.

For a diner, a review-based search engine offers:
\begin{itemize}
  \item search in their own words instead of browsing categories and filters;
  \item a reason for every recommendation, in the form of real review sentences they can read;
  \item better handling of specific constraints such as price, waiting time, atmosphere and,
  in Thai, taste words such as \thai{เผ็ด} (spicy) or \thai{ไม่หวานมาก} (not too sweet).
\end{itemize}

For a restaurant platform, the same engine can surface restaurants whose strengths appear in
reviews but not in category labels, and it reduces the trust gap caused by opaque rankings.
Thailand is a natural market for this idea: Wongnai alone contains tens of thousands of Thai
reviews, and Thai diners often describe food in terms of taste. The system does not replace
the user's judgement. Reviews can be subjective, outdated or contradictory, and the system
presents them as evidence to read, not as verified facts.

\section{Stakeholders and success criteria}
The primary stakeholder is the diner, who needs relevant choices and evidence they can
understand. The secondary stakeholder is the platform operator, who needs a service that is
measurable, affordable to run and robust to the uneven number of reviews that restaurants
receive. We define success with four criteria:
\begin{enumerate}
  \item \textbf{Ranking quality:} the proposed model ranks restaurants better than a strong
  lexical baseline (BM25) under a fair comparison.
  \item \textbf{Evidence quality:} the shown sentences come from the original reviews, are
  traceable to their source, and are tested for whether they drive and support the ranking.
  \item \textbf{Language transfer:} the same retrieval method is tested on real Thai reviews
  with human relevance labels created for this project, and the result is reported whether or
  not it is favourable.
  \item \textbf{Practicality:} the system is simple enough to run in an interactive demo.
\end{enumerate}

\section{Research objectives}
The project has three objectives, one for each part of the guiding question:
\begin{enumerate}
  \item compare BM25 and multilingual dense retrieval fairly, using the same text units and
  the same aggregation functions, and measure how review volume affects the result;
  \item evaluate whether the shown review sentences are faithful to the ranking and support the
  query;
  \item test whether sentence retrieval transfers to real Thai restaurant reviews, using taste
  queries that were fixed in advance and labelled by a native Thai speaker.
\end{enumerate}

\section{Contributions}
The main contributions of this project are:
\begin{enumerate}
  \item \textbf{A fair comparison on RIRD.} BM25 and a multilingual dense encoder score the
  same units (reviews or sentences) with the same aggregation functions, with bootstrap
  confidence intervals and paired significance tests. The proposed configuration, dense
  sentence retrieval with top-10 mean aggregation, reaches nDCG@10 0.682 and MAP 0.618 without
  fine-tuning, close to the fine-tuned RIRD reference (MAP 0.626).
  \item \textbf{A review-volume control.} In the review-level experiments, capping every
  restaurant at 400 reviews removes the correlation between restaurant score and review count,
  and changes nDCG@10 by at most 0.011.
  \item \textbf{An explanation evaluation protocol.} The shown sentences are tested with a
  deletion test for faithfulness and a human rating for support.
  \item \textbf{A Thai taste test set.} To our knowledge, this is the first small
  human-labelled test set for taste-based retrieval over Thai restaurant reviews: 24 frozen
  queries and 452 labelled query--sentence pairs. On this set, BM25 (nDCG@10 0.654) scores
  higher than the dense model (0.549), but the difference is not statistically significant.
  This is the opposite of the English result.
  \item \textbf{A working demo.} An interactive application shows ranked restaurants with
  their evidence sentences in English and labelled sentence results in Thai, with query
  rewriting and selective cross-encoder reranking for queries that contain negation.
\end{enumerate}

\section{Organisation of the report}
Chapter~\ref{ch:problem} defines the task formally and states the research questions.
Chapter~\ref{ch:related} reviews related work and positions the project. Chapter~\ref{ch:data}
describes the two datasets and the construction of the Thai test set.
Chapter~5 presents the methodology, Chapter~6 the evaluation results, Chapter~7 the
discussion, and Chapter~8 the conclusions.
